\chapter{How D'Artagnan Became Acquainted  with a Poet}

\lettrine{B}{efore} taking his place at table, D'Artagnan  acquired, as was his custom, all the information he could; but it is an axiom  of curiosity, that every man who wishes to question well and fruitfully ought  in the first place to lay himself open to questions. D'Artagnan sought,  then, with his usual skill, a promising questioner in the hostelry of La  Roche-Bernard. At the moment, there were in the house, on the first story, two  travelers either preparing for supper, or at supper itself. D'Artagnan  had seen their nags in the stable, and their equipages in the salle. One  traveled with a lackey, undoubtedly a person of consideration;—two Perche  mares, sleek, sound beasts, were suitable means of locomotion. The other, a  little fellow, a traveller of meagre appearance, wearing a dusty surtout, dirty  linen, and boots more worn by the pavement than the stirrup, had come from  Nantes with a cart drawn by a horse so like Furet in colour, that  D'Artagnan might have gone a hundred miles without finding a better  match. This cart contained divers large packets wrapped in pieces of old stuff.

<That traveller yonder,> said D'Artagnan to himself, <is  the man for my money. He will do, he suits me; I ought to do for him and suit  him; M. Agnan, with the gray doublet and the rusty calotte, is not unworthy of  supping with the gentleman of the old boots and still older horse.>

This said, D'Artagnan called the host, and desired him to send his teal,  tourteau, and cider up to the chamber of the gentleman of modest exterior. He  himself climbed, a plate in his hand, the wooden staircase which led to the  chamber, and began to knock at the door.

<Come in!> said the unknown. D'Artagnan entered, with a  simper on his lips, his plate under his arm, his hat in one hand, his candle in  the other.

<Excuse me, monsieur,> said he, <I am as you are, a traveller;  I know no one in the hotel, and I have the bad habit of losing my spirits when  I eat alone; so that my repast appears a bad one to me, and does not nourish  me. Your face, which I saw just now, when you came down to have some oysters  opened,—your face pleased me much. Besides, I have observed you have a  horse just like mine, and that the host, no doubt on account of that  resemblance, has placed them side by side in the stable, where they appear to  agree amazingly well together. I therefore, monsieur, do not see any reason why  the masters should be separated when the horses are united. Accordingly, I am  come to request the pleasure of being admitted to your table. My name is Agnan,  at your service, monsieur, the unworthy steward of a rich seigneur, who wishes  to purchase some salt-mines in this country, and sends me to examine his future  acquisitions. In truth, monsieur, I should be well pleased if my countenance  were as agreeable to you as yours is to me; for, upon my honour, I am quite at  your service.>

The stranger, whom D'Artagnan saw for the first time,—for before he  had only caught a glimpse of him,—the stranger had black and brilliant  eyes, a yellow complexion, a brow a little wrinkled by the weight of fifty  years, bonhomie in his features collectively, but some cunning in his look.

<One would say,> thought D'Artagnan, <that this merry  fellow has never exercised more than the upper part of his head, his eyes, and  his brain. He must be a man of science: his mouth, nose, and chin signify  absolutely nothing.>

<Monsieur,> replied the latter, with whose mind and person we have  been making so free, <you do me much honour; not that I am ever ennuye,  for I have,> added he, smiling, <a company which amuses me always:  but, never mind that, I am happy to receive you.> But when saying this,  the man with the worn boots cast an uneasy look at his table, from which the  oysters had disappeared, and upon which there was nothing left but a morsel of  salt bacon.

<Monsieur,> D'Artagnan hastened to say, <the host is  bringing me up a pretty piece of roasted poultry and a superb tourteau.>  D'Artagnan had read in the look of his companion, however rapidly it  disappeared, the fear of an attack by a parasite: he divined justly. At this  opening, the features of the man of modest exterior relaxed; and, as if he had  watched the moment for his entrance, as D'Artagnan spoke, the host  appeared, bearing the announced dishes. The tourteau and the teal were added to  the morsel of broiled bacon; D'Artagnan and his guest bowed, sat down  opposite to each other, and, like two brothers, shared the bacon and the other  dishes.

<Monsieur,> said D'Artagnan, <you must confess that  association is a wonderful thing.>

<How so?> replied the stranger, with his mouth full.

<Well, I will tell you,> replied D'Artagnan.

The stranger gave a short truce to the movement of his jaws, in order to hear  the better.

<In the first place,> continued D'Artagnan, <instead of  one candle, which each of us had, we have two.>

<That is true!> said the stranger, struck with the extreme lucidity  of the observation.

<Then I see that you eat my tourteau in preference, whilst I, in  preference, eat your bacon.>

<That is true again.>

<And then, in addition to being better lighted and eating what we prefer,  I place the pleasure of your company.>

<Truly, monsieur, you are very jovial,> said the unknown,  cheerfully.

<Yes, monsieur; jovial, as all people are who carry nothing on their  minds, or, for that matter, in their heads. Oh! I can see it is quite another  sort of thing with you,> continued D'Artagnan; <I can read in  your eyes all sorts of genius.>

<Oh, monsieur!>

<Come, confess one thing.>

<What is that?>

<That you are a learned man.>

<Ma foi! monsieur.>

<Hein?>

<Almost.>

<Come, then!>

<I am an author.>

<There!> cried D'Artagnan, clapping his hands, <I knew  I could not be deceived! It is a miracle!>

<Monsieur\longdash>

<What, shall I have the honour of passing the evening in the society of an  author, of a celebrated author, perhaps?>

<Oh!> said the unknown, blushing, <celebrated, monsieur,  celebrated is not the word.>

<Modest!> cried D'Artagnan, transported, <he is  modest!> Then, turning towards the stranger, with a character of blunt  bonhomie: <But tell me at least the name of your works, monsieur; for you  will please to observe you have not told me your name, and I have been forced  to divine your genius.>

<My name is Jupenet, monsieur,> said the author.

<A fine name! a grand name! upon my honour; and I do not know  why—pardon me the mistake, if it be one—but surely I have heard  that name somewhere.>

<I have made verses,> said the poet, modestly.

<Ah! that is it, then; I have heard them read.>

<A tragedy.>

<I must have seen it played.>

The poet blushed again, and said: <I do not think that can be the case,  for my verses have never been printed.>

<Well, then, it must have been the tragedy which informed me of your  name.>

<You are again mistaken, for MM. the comedians of the Hotel de Bourgogne,  would have nothing to do with it,> said the poet, with a smile, the  receipt for which certain sorts of pride alone knew the secret.  D'Artagnan bit his lips. <Thus, then, you see, monsieur,>  continued the poet, <you are in error on my account, and that not being  at all known to you, you have never heard tell of me.>

<Ah! that confounds me. That name, Jupenet, appears to me, nevertheless,  a fine name, and quite as worthy of being known as those of MM. Corneille, or  Rotrou, or Garnier. I hope, monsieur, you will have the goodness to repeat to  me a part of your tragedy presently, by way of dessert, for instance. That will  be sugared roast meat,—mordioux! Ah! pardon me, monsieur, that was a  little oath which escaped me, because it is a habit with my lord and master. I  sometimes allow myself to usurp that little oath, as it seems in pretty good  taste. I take this liberty only in his absence, please to observe, for you may  understand that in his presence—but, in truth, monsieur, this cider is  abominable; do you not think so? And besides, the pot is of such an irregular  shape it will not stand on the table.>

<Suppose we were to make it level?>

<To be sure; but with what?>

<With this knife.>

<And the teal, with what shall we cut that up? Do you not, by chance,  mean to touch the teal?>

<Certainly.>

<Well, then\longdash>

<Wait.>

And the poet rummaged in his pocket, and drew out a piece of brass, oblong,  quadrangular, about a line in thickness, and an inch and a half in length. But  scarcely had this little piece of brass seen the light, than the poet appeared  to have committed an imprudence, and made a movement to put it back again in  his pocket. D'Artagnan perceived this, for he was a man that nothing  escaped. He stretched forth his hand towards the piece of brass: <Humph!  that which you hold in your hand is pretty; will you allow me to look at  it?>

<Certainly,> said the poet, who appeared to have yielded too soon  to a first impulse. <Certainly, you may look at it: but it will be in  vain for you to look at it,> added he, with a satisfied air; <if I  were not to tell you its use, you would never guess it.>

D'Artagnan had seized as an avowal the hesitation of the poet, and his  eagerness to conceal the piece of brass which a first movement had induced him  to take out of his pocket. His attention, therefore, once awakened on this  point, he surrounded himself with a circumspection which gave him a superiority  on all occasions. Besides, whatever M. Jupenet might say about it, by a simple  inspection of the object, he perfectly well knew what it was. It was a  character in printing.

<Can you guess, now, what this is?> continued the poet.

<No,> said D'Artagnan, <no, ma foi!>

<Well, monsieur,> said M. Jupenet, <this little piece of  metal is a printing letter.>

<Bah!>

<A capital.>

<Stop, stop, stop,> said D'Artagnan, opening his eyes very  innocently.

<Yes, monsieur, a capital; the first letter of my name.>

<And this is a letter, is it?>

<Yes, monsieur.>

<Well, I will confess one thing to you.>

<And what is that?>

<No, I will not, I was going to say something stupid.>

<No, no,> said Master Jupenet, with a patronizing air.

<Well, then, I cannot comprehend, if that is a letter, how you can make a  word.>

<A word?>

<Yes, a printed word.>

<Oh, that's very easy.>

<Let me see.>

<Does it interest you?>

<Enormously.>

<Well, I will explain the thing to you. Attend.>

<I am attending.>

<This is it.>

<Good.>

<Look attentively.>

<I am looking.> D'Artagnan, in fact, appeared absorbed in  observations. Jupenet drew from his pocket seven or eight other pieces of brass  smaller than the first.

<Ah, ah,> said D'Artagnan.

<What!>

<You have, then, a whole printing-office in your pocket. Peste! that is  curious, indeed.>

<Is it not?>

<Good God, what a number of things we learn by traveling.>

<To your health!> said Jupenet, quite enchanted.

<To yours, mordioux, to yours. But—an instant—not in this  cider. It is an abominable drink, unworthy of a man who quenches his thirst at  the Hippocrene fountain—is not it so you call your fountain, you  poets?>

<Yes, monsieur, our fountain is so called. That comes from two Greek  words—hippos, which means a horse, and\longdash>

<Monsieur,> interrupted D'Artagnan, <you shall drink of  a liquor which comes from one single French word, and is none the worse for  that—from the word grape; this cider gives me the heartburn. Allow me to  inquire of your host if there is not a good bottle of Beaugency, or of the  Ceran growth, at the back of the large bins in his cellar.>

The host, being sent for, immediately attended.

<Monsieur,> interrupted the poet, <take care, we shall not  have time to drink the wine, unless we make great haste, for I must take  advantage of the tide to secure the boat.>

<What boat?> asked D'Artagnan.

<Why the boat which sets out for Belle-Isle.>

<Ah—for Belle-Isle,> said the musketeer, <that is  good.>

<Bah! you will have plenty of time, monsieur,> replied the  hotelier, uncorking the bottle, <the boat will not leave this  hour.>

<But who will give me notice?> said the poet.

<Your fellow-traveller,> replied the host.

<But I scarcely know him.>

<When you hear him departing, it will be time for you to go.>

<Is he going to Belle-Isle, likewise, then?>

<Yes.>

<The traveller who has a lackey?> asked D'Artagnan. <He  is some gentleman, no doubt?>

<I know nothing of him.>

<What!—know nothing of him?>

<No, all I know is, that he is drinking the same wine as you.>

<Peste!—that is a great honour for us,> said D'Artagnan,  filling his companion's glass, whilst the host went out.

<So,> resumed the poet, returning to his dominant ideas, <you  never saw any printing done?>

<Never.>

<Well, then, take the letters thus, which compose the word, you see: A B;  ma foi! here is an R, two E E, then a G.> And he assembled the letters  with a swiftness and skill which did not escape the eye of D'Artagnan.

<Abrege,> said he, as he ended.

<Good!> said D'Artagnan; <here are plenty of letters  got together; but how are they kept so?> And he poured out a second glass  for the poet. M. Jupenet smiled like a man who has an answer for everything;  then he pulled out—still from his pocket—a little metal ruler,  composed of two parts, like a carpenter's rule, against which he put  together, and in a line, the characters, holding them under his left thumb.

<And what do you call that little metal ruler?> said  D'Artagnan, <for, I suppose, all these things have names.>

<This is called a composing-stick,> said Jupenet; <it is by  the aid of this stick that the lines are formed.>

<Come, then, I was not mistaken in what I said; you have a press in your  pocket,> said D'Artagnan, laughing with an air of simplicity so  stupid, that the poet was completely his dupe.

<No,> replied he; <but I am too lazy to write, and when I  have a verse in my head, I print it immediately. That is a labour spared.>

<Mordioux!> thought D'Artagnan to himself, <this must  be cleared up.> And under a pretext, which did not embarrass the  musketeer, who was fertile in expedients, he left the table, went downstairs,  ran to the shed under which stood the poet's little cart, and poked the  point of his poniard into the stuff which enveloped one of the packages, which  he found full of types, like those which the poet had in his pocket.

<Humph!> said D'Artagnan, <I do not yet know whether M.  Fouquet wishes to fortify Belle-Isle; but, at all events, here are some  spiritual munitions for the castle.> Then, enchanted with his rich  discovery, he ran upstairs again, and resumed his place at the table.

D'Artagnan had learnt what he wished to know. He, however, remained, none  the less, face to face with his partner, to the moment when they heard from the  next room symptoms of a person's being about to go out. The printer was  immediately on foot; he had given orders for his horse to be got ready. His  carriage was waiting at the door. The second traveller got into his saddle, in  the courtyard, with his lackey. D'Artagnan followed Jupenet to the door;  he embarked his cart and horse on board the boat. As to the opulent traveller,  he did the same with his two horses and servant. But all the wit  D'Artagnan employed in endeavouring to find out his name was lost—he  could learn nothing. Only he took such notice of his countenance, that it was  impressed upon his mind forever. D'Artagnan had a great inclination to  embark with the two travelers, but an interest more powerful than  curiosity—that of success—repelled him from the shore, and brought  him back again to the hostelry. He entered with a sigh, and went to bed  directly in order to be ready early in the morning with fresh ideas and the  sage counsel of sufficing sleep.  